Podíl kolektivních smluv obsahujících formální mzdovou tarifní soustavu, závazek sociálního klidu a~ustanovení o~vzdělávání dokumentuje, že česká \ac{KS} jsou obsahově chudší, než jejich počet naznačuje.
Přítomnost tarifní mzdové soustavy --- systému platových tříd a~tarifů, který zabraňuje individuálnímu sjednávání pod tlakem --- klesá, přičemž \ac{KS} stále více fungují jako formální stvrzení zákonného minima bez skutečné koordinační přidané hodnoty.
Tento „hollowing out“ obsahu \ac{KS} je symptomatický pro fragmentaci na podnikové úrovni: bez odvětvového referenčního rámce nemají zástupci zaměstnanců ani zaměstnavatelů motivaci investovat do sofistikovaných mzdových struktur, jejichž přínos by byl viditelný pouze v~systémovém srovnání.
Reforma zaměřená na rozšíření závaznosti odvětvových \ac{KS} (erga omnes) by přirozeně přinesla i~kvalitnější obsah smluv, protože sektorové smlouvy musejí pokrývat diverzní podnikové situace a~tudíž systémovost tarifních struktur přímo vyžadují.

Podíl kolektivních smluv obsahujících formální mzdovou tarifní soustavu, závazek sociálního klidu a~ustanovení o~vzdělávání dokumentuje, že česká \ac{KS} jsou obsahově chudší, než jejich počet naznačuje.
Přítomnost tarifní mzdové soustavy --- systému platových tříd a~tarifů, který zabraňuje individuálnímu sjednávání pod tlakem --- klesá, přičemž \ac{KS} stále více fungují jako formální stvrzení zákonného minima bez skutečné koordinační přidané hodnoty.
Tento „hollowing out“ obsahu \ac{KS} je symptomatický pro fragmentaci na podnikové úrovni: bez odvětvového referenčního rámce nemají zástupci zaměstnanců ani zaměstnavatelů motivaci investovat do sofistikovaných mzdových struktur, jejichž přínos by byl viditelný pouze v~systémovém srovnání.
Reforma zaměřená na rozšíření závaznosti odvětvových \ac{KS} (erga omnes) by přirozeně přinesla i~kvalitnější obsah smluv, protože sektorové smlouvy musejí pokrývat diverzní podnikové situace a~tudíž systémovost tarifních struktur přímo vyžadují.

Podíl kolektivních smluv obsahujících formální mzdovou tarifní soustavu, závazek sociálního klidu a~ustanovení o~vzdělávání dokumentuje, že česká \ac{KS} jsou obsahově chudší, než jejich počet naznačuje.
Přítomnost tarifní mzdové soustavy --- systému platových tříd a~tarifů, který zabraňuje individuálnímu sjednávání pod tlakem --- klesá, přičemž \ac{KS} stále více fungují jako formální stvrzení zákonného minima bez skutečné koordinační přidané hodnoty.
Tento „hollowing out“ obsahu \ac{KS} je symptomatický pro fragmentaci na podnikové úrovni: bez odvětvového referenčního rámce nemají zástupci zaměstnanců ani zaměstnavatelů motivaci investovat do sofistikovaných mzdových struktur, jejichž přínos by byl viditelný pouze v~systémovém srovnání.
Reforma zaměřená na rozšíření závaznosti odvětvových \ac{KS} (erga omnes) by přirozeně přinesla i~kvalitnější obsah smluv, protože sektorové smlouvy musejí pokrývat diverzní podnikové situace a~tudíž systémovost tarifních struktur přímo vyžadují.

Podíl kolektivních smluv obsahujících formální mzdovou tarifní soustavu, závazek sociálního klidu a~ustanovení o~vzdělávání dokumentuje, že česká \ac{KS} jsou obsahově chudší, než jejich počet naznačuje.
Přítomnost tarifní mzdové soustavy --- systému platových tříd a~tarifů, který zabraňuje individuálnímu sjednávání pod tlakem --- klesá, přičemž \ac{KS} stále více fungují jako formální stvrzení zákonného minima bez skutečné koordinační přidané hodnoty.
Tento „hollowing out“ obsahu \ac{KS} je symptomatický pro fragmentaci na podnikové úrovni: bez odvětvového referenčního rámce nemají zástupci zaměstnanců ani zaměstnavatelů motivaci investovat do sofistikovaných mzdových struktur, jejichž přínos by byl viditelný pouze v~systémovém srovnání.
Reforma zaměřená na rozšíření závaznosti odvětvových \ac{KS} (erga omnes) by přirozeně přinesla i~kvalitnější obsah smluv, protože sektorové smlouvy musejí pokrývat diverzní podnikové situace a~tudíž systémovost tarifních struktur přímo vyžadují.

\input{texparts/python/ipp_ca_breadth}



